\subsection{Estrategia de verificación y validación}

La verificación comprueba si el producto satisface los requisitos especificados; la validación determina si la solución responde a las necesidades del personal asistencial en su contexto de uso. En Hito 3 ambas actividades tienen un alcance preliminar: se realizó una inspección del código y de los flujos disponibles, pero no existe una suite automatizada ni una prueba formal con usuarios clínicos.

Para evitar conclusiones no sustentadas, se utilizan tres resultados de verificación:

\begin{itemize}
    \item \textbf{Cumplido:} existe evidencia funcional suficiente para el alcance definido.
    \item \textbf{Parcial:} existe una implementación, pero faltan condiciones esenciales o pruebas.
    \item \textbf{Pendiente:} no existe una implementación reproducible en el prototipo actual.
\end{itemize}

La evaluación se efectuó mediante trazabilidad entre requisitos, rutas de interfaz, endpoints, modelos de datos y servicios disponibles. Este método permite reconocer brechas, pero no sustituye pruebas unitarias, de integración, seguridad, rendimiento y usabilidad.
